\chapter{Модулярный двойственный вес и коэрцитивность мостовой долины}

Пространственная жёсткость не поднимает однородную долину. Остаётся
проверить структурный намёк Тома V: вещественная спектральная геометрия
содержит левое действие, противоположное действие и модулярный оператор
\begin{equation}
 \Delta_R(X)=RXR^{-1}.
\end{equation}
Может ли эта тройка автоматически оценивать не только опору $R$, но и её
ядро?

\section{Прямое и противоположное действия}

Для вещественного симметричного семейного состояния
\begin{equation}
 R=R^T
\end{equation}
вещественная структура переводит левый вес в транспонированный правый.
Для мостового функционала это даёт тождество
\begin{equation}
 \mathcal S_R(B)=\mathcal S_{R^T}(B^T).
 \label{eq:v6-modular-dual-transpose-identity}
\end{equation}
Следовательно, сумма прямой и коммутантной копий равна
$2\mathcal S_R(B)$ и не создаёт нового дополнения состояния.

На проекторной страте
\begin{equation}
 R=P=\operatorname{diag}(1,0,0),
 \qquad B_C=\operatorname{diag}(1,C)
\end{equation}
обе копии равны нулю для любого $C\in M_2(\mathbb R)$. Противоположное
действие содержит $R^T=R$, а не $I-R$.

\begin{proposition}[Коммутант не является дополнением состояния]
Real-сопряжение сохраняет спектр и ранг семейного состояния. Оно не
переводит опору $R$ в проектор на $\ker R$, поэтому само по себе не
устраняет потерю коэрцитивности.
\end{proposition}

\section{КМС-симметричная норма}

Естественная симметризация скалярного произведения имеет вид
\begin{equation}
 \langle X,X\rangle_{\rm KMS}
 =\Tr\bigl(R^{1/2}X^*R^{1/2}X\bigr).
 \label{eq:v6-modular-kms-norm}
\end{equation}
Более общие монотонные квантовые метрики и среднее Боголюбова--Кубо--Мори
также строятся на многообразии верных состояний; см.
\path{arXiv:1008.2417}, \path{arXiv:0909.3647} и
\path{arXiv:2203.10857}.

Возьмём
\begin{equation}
 R_\varepsilon
 =\operatorname{diag}(1-2\varepsilon,\varepsilon,\varepsilon).
\end{equation}
Для дефекта, целиком лежащего в поперечном блоке $C$, как GNS-норма, так
и КМС-норма пропорциональны $\varepsilon$ и стремятся к нулю. На
$\varepsilon=0$ они снова не видят поперечный блок.

Машинный аудит подтвердил это на случайных $C$: прямой, коммутантный и
КМС-функционалы точно равны нулю на проекторной долине.

\section{Обратный модулярный вес}

Формальный функционал
\begin{equation}
 \Tr\bigl(R^{-1}X^*X\bigr)
 \label{eq:v6-modular-inverse-candidate}
\end{equation}
действительно расходится как $\varepsilon^{-1}$ и подавляет ядро. Но из
существования модулярного оператора не следует, что именно
\eqref{eq:v6-modular-inverse-candidate} входит в энергию.

В Томe V состояние
\begin{equation}
 \rho_\beta=\frac{e^{-\beta h_F}}{\Tr e^{-\beta h_F}}
\end{equation}
использовалось для поляризации частот и определения ориентированного
дифференциала. Оно верно при конечном $\beta$, а параметр $\beta$ не был
выведен как наблюдаемое число. Квадрат кривизны нормировался обычным
матричным следом, не обратным состоянием. Кроме того, семейное $R$ Тома VI
не было отождествлено с $\rho_\beta$.

\begin{theorem}[Обратный вес не следует из модулярной кинематики]
Оператор $L_RR_{R^{-1}}$ канонически существует для верного состояния и
становится неограниченным у границы ранга. Но текущая архитектура выводит
из него ориентацию потока, а не положительный энергетический член
$\Tr(R^{-1}X^*X)$. Добавление такого члена было бы новой модельной
аксиомой.
\end{theorem}

\section{Каноническая относительная энтропия}

Можно избежать выбора произвольного оператора, сравнивая $R$ с
трассовым состоянием $R_0=I_3/3$. Обратная относительная энтропия равна
\begin{equation}
 D(R_0\|R)
 =\Tr R_0(\log R_0-\log R)
 =-\frac13\log\det R-\log3.
 \label{eq:v6-modular-relative-entropy-barrier}
\end{equation}
Она действительно расходится при потере ранга и не требует выбора оси.
Однако её коэффициент барьера равен $1/3$, что значительно превышает
$17/168$. Вблизи $I_3/3$ она добавляет кривизну $3/2$, и общий коэффициент
становится
\begin{equation}
 -\frac{51}{112}+\frac32
 =\frac{117}{112}>0.
\end{equation}
Самозапуск исчезает.

Относительная энтропия не обязана входить в действие только потому, что
два состояния доступны. Но даже наиболее канонический выбор не решает
задачу: он стабилизирует смешанный вакуум слишком сильно.

\section{Вердикт}

\begin{protocol}
\begin{enumerate}
 \item прямая и коммутантная мостовые нормы:
 \textbf{совпадают};
 \item противоположный вес равен $I-R$:
 \textbf{нет};
 \item проекторная долина после Real-удвоения:
 \textbf{сохраняется};
 \item КМС-симметричная норма поднимает $\ker R$:
 \textbf{нет};
 \item обратный вес $R^{-1}$ создаёт барьер:
 \textbf{формально да};
 \item обратный вес выведен родительским действием:
 \textbf{нет};
 \item трассовая относительная энтропия:
 \textbf{даёт барьер $1/3$};
 \item отрицательная мода после этого барьера:
 \textbf{исчезает};
 \item модулярная двойственная коэрцитивность:
 \textbf{не получена};
 \item следующая проверка:
 \textbf{самосогласованная очистка состояния мостом};
 \item рождение материи:
 \textbf{не закрыто}.
\end{enumerate}
\end{protocol}

Следующий тест должен снять независимость $R$ и $B$. Если семейное
состояние является редуцированным состоянием или очищением самого моста,
то условие $R=P$ может принудить $C=0$ и устранить долину без нового
потенциала. Это должно быть проверено на уже построенной qutrit-purification
кинематике, а не введено как словарь после результата.

Машинный сертификат сохранён в
\path{s2t_v6_modular_dual_weight_bridge_coercivity_gate_results.json}.